Put \(L=\log(en)\), \(u=d^2/(n^2L^2)\), and \(q=d/n^2\).  At \(\sigma=0\), the definition gives
\[
 r_{n,d,0}=n^{-1}+\min\{1,u,q\}.
\tag{1}
\]
The upper comparison \(r_{n,d,0}\le n^{-1}+q\) is immediate.  If \(d\ge L^2\), then \(u\ge q\); the stipulated range, after taking \(c_{\epsilon}\le1\), also gives \(q\le1\), so the minimum in (1) equals \(q\).  If \(d<L^2\), then \(q<L^2/n^2\le C/n\), because \(\sup_{n\ge1}L^2/n<\infty\).  Hence in both cases
\[
 r_{n,d,0}\asymp n^{-1}+q.
\tag{2}
\]
The lower expression in the revised theorem is exactly \(n^{-1}+q\) at \(\sigma=0\), so (2) proves matching at the first endpoint.

At \(\sigma=2\), \(4+q\ge1\), and therefore
\[
 r_{n,d,2}=n^{-1}+\min\{1,u\}.
\tag{3}
\]
Lemma~\(\mathrm{lem:scale\text{-}sanity}\) gives \(\mathcal P_{d,\epsilon,M,2}=\mathcal P_{d,\epsilon,M}^{\mathrm{unr}}\).  The lower expression at this endpoint is \(n^{-1}+q+4\min\{1,u\}\).  It remains only to check that \(q\) is dominated by (3).  If \(d\ge L^2\), then \(q\le u\); if \(d<L^2\), then \(q\le C/n\), exactly as above.  Thus
\[
 n^{-1}+q+4\min\{1,u\}\asymp n^{-1}+\min\{1,u\}.
\tag{4}
\]
Equations (3)--(4) and the class equality prove matching at the unrestricted endpoint.  All comparisons are uniform on the stated dimension range, including zero-mass cells and the boundary values of \(\sigma\).